% BHU PYQ -- Manually transcribed & error-corrected
% Paper: GLB-502: Igneous Petrology, Mineralogy and Crystallography
% Exam: B.Sc. (Hons.) Semester V Examination, 2016-17

\documentclass[11pt,a4paper]{article}
\usepackage[a4paper,top=2.5cm,bottom=2.5cm,left=2.8cm,right=2.8cm,headheight=14pt]{geometry}
\usepackage{amsmath,amssymb}
\usepackage[T1]{fontenc}
\usepackage[utf8]{inputenc}
\usepackage{lmodern,microtype,fancyhdr,enumitem,hyperref,lastpage,parskip}

\hypersetup{colorlinks=true,urlcolor=black,linkcolor=black,
  pdftitle={GLB-502 Igneous Petrology, Mineralogy and Crystallography -- BHU B.Sc. Sem V 2016-17}}
\pagestyle{fancy}\fancyhf{}
\renewcommand{\headrulewidth}{0.4pt}\renewcommand{\footrulewidth}{0.4pt}
\fancyhead[L]{\small   \textit{Banaras Hindu University}}
\fancyhead[R]{\small   \textit{GLB-502: Igneous Petrology, Mineralogy \& Crystallography}}
\fancyfoot[C]{\small Page  \thepage\ of \pageref{LastPage}}


\newlist{parts}{enumerate}{2}
\setlist[parts,1]{label=(\alph*),leftmargin=*,itemsep=5pt,topsep=3pt}

\newlist{romanparts}{enumerate}{2}
\setlist[romanparts,1]{label=(\roman*),leftmargin=*,itemsep=5pt,topsep=3pt}

\newcommand{\pts}[1]{\hfill   \textit{[#1]}}


\begin{document}

\begin{center}
  {\large\bfseries BANARAS HINDU UNIVERSITY}\\[2pt]
  {
\normalsize B.Sc. (Hons.) --- Semester V Examination, 2016-17}\\[6pt]
  \rule{0.8\linewidth}{0.5pt}\\[6pt]
  {\large\bfseries Paper No. GLB-502: Igneous Petrology, Mineralogy and Crystallography}\\[4pt]
  {
\normalsize Time: 3 Hours \hfill Maximum Marks: 70}\\[2pt]
  {\small Ref. No.: BS/Sem V/73}
\end{center}
\rule{\linewidth}{0.4pt}
\smallskip



\noindent   \textit{Instructions: Attempt any   \textbf{five} questions, selecting at least   \textbf{two} from each of Sections A and B. Section C is compulsory. All questions carry equal marks (14 marks each).}
\bigskip

\bigskip


\noindent {\large\bfseries SECTION --- A}
\rule{\linewidth}{0.4pt}\smallskip




\noindent   \textbf{Question 1.} Give the salient features of chemical classification schemes for igneous rocks. \pts{14}

\medskip


\noindent   \textbf{Question 2.} Write the petrographic description of   \textbf{any four} of the following: \pts{$4  \times 3.5 = 14$}

\begin{parts}
  \item Norite
  \item Komatiite
  \item Rhyolite
  \item Andesite
  \item Dunite
  \item Phonolite
\end{parts}

\medskip


\noindent   \textbf{Question 3.} Write descriptive notes on the following: \pts{$2  \times 7 = 14$}

\begin{parts}
  \item Fractional crystallization
  \item Types of magmas
\end{parts}

\medskip


\noindent   \textbf{Question 4.} Describe with the help of a neat sketch the crystallization of magma in the phase equilibrium of Diopside-Albite-Anorthite system and comment on its petrogenetic significance. \pts{14}

\bigskip


\noindent {\large\bfseries SECTION --- B}
\rule{\linewidth}{0.4pt}\smallskip




\noindent   \textbf{Question 5.} Describe the symmetry and forms belonging to the holosymmetric class of the Monoclinic system and draw a suitable stereogram. Name any four minerals belonging to this class. \pts{14}

\medskip


\noindent   \textbf{Question 6.} Write short notes on   \textbf{any two} of the following: \pts{$2  \times 7 = 14$}

\begin{parts}
  \item Optic sign
  \item Solid solution
  \item Plagioclase feldspars
\end{parts}

\medskip


\noindent   \textbf{Question 7.} Discuss in detail the structural, physical, chemical, and optical properties of the mica group of minerals. \pts{14}

\bigskip


\noindent {\large\bfseries SECTION --- C (Compulsory)}
\rule{\linewidth}{0.4pt}\smallskip




\noindent   \textbf{Question 8.} Answer the following questions (each carry 2 marks): \pts{$7  \times 2 = 14$}

\begin{romanparts}
  \item Name any two igneous rocks which are ultramafic as well as ultrabasic.
  \item What are the factors that control the stability of minerals?
  \item What is the utility of the Michel-Lévy Chart?
  \item How do we know that there is water in the Earth's mantle?
  \item What is the difference between eutectic and cotectic?
  \item How many optic axes are there in minerals crystallizing in the Triclinic system?
  \item What is the extinction angle in quartz?
\end{romanparts}

\vfill\rule{\linewidth}{0.4pt}

\begin{center}\small
\href{https://bhu-exam.vercel.app/}{Click here to visit our website}\\[4pt]
\href{https://me-aryan.vercel.app/}{Click here to visit our contributor}\\[4pt]
\href{https://quashan-me.vercel.app/}{Click here to visit our contributor}
\end{center}
\end{document}
